This chapter explains the integration and communication layout of the AgriVerify platform. The application is divided into three distinct areas: a Next.js web frontend, an ASP.NET Core backend server for business logic, and a FastAPI service that serves our machine learning models \cite{martin2018clean}. Decoupling the frontend user interface, data validation, and deep learning inference allows each service to run and scale independently based on demand.

\section{System Integration}

The AgriVerify system routes communication across several different boundaries. During early development, hosting the Next.js frontend and ASP.NET Core backend on separate subdomains caused CORS blocking in browsers. To solve this, we implemented a dedicated API proxy route in Next.js \cite{nextjs2024}. The browser sends requests directly to the local Next.js server, which intercepts the call, injects authentication tokens, and forwards the payload to the internal ASP.NET Core backend. This prevents exposing internal API paths and security credentials directly to the public web \cite{masue2024quantifying}. Table~\ref{tab:integration_matrix} lists the communication paths, protocols, and security methods used across the services.

\begin{table}[htbp]
\centering
\caption{System Integration and Communication Matrix}
\label{tab:integration_matrix}
\small
\renewcommand{\arraystretch}{1.1}
\begin{tabularx}{\textwidth}{|p{3.5cm}|p{3.5cm}|X|}
\hline
\textbf{Communication Path} & \textbf{Protocol} & \textbf{Technical Implementation \& Security} \\
\hline
\textbf{Browser to Proxy} & HTTPS with multipart uploads & Client-side Axios interceptors automatically inject JSON Web Tokens (JWT) into request headers and coordinate silent token rotation upon expiry \cite{axios2024, jones2015jwt}. \\
\hline
\textbf{Proxy to Backend} & REST API over HTTPS & The proxy gateway intercepts incoming requests, performs schema validation and forwards payloads to internal ASP.NET Core endpoints, encapsulating internal network topology \cite{freeman2020pro}. \\
\hline
\textbf{Backend to AI} & REST (HTTPS) & The backend orchestrates parallel requests, transmitting image binaries to dedicated AI services for quality checks (blur/illumination detection) and Quality Assessment models (ResNet50) \cite{he2016deep}. \\
\hline
\textbf{Cloud Services} & Azure SDK Client Libraries & Azure OpenAI Service for text synthesis and Azure Blob Storage for unstructured binary object persistence \cite{copard2020azure}. \\
\hline
\textbf{Deployment} & Container Orchestration & Multi-stage Docker builds compile and optimize service codebases, which are stored in the Azure Container Registry and deployed using container orchestration tools \cite{merkel2014docker}. \\
\hline
\end{tabularx}
\end{table}

\section{End-to-End Architectural Workflow}

Handling image files from modern mobile devices presented a significant design challenge. High-resolution cameras on mid-range smartphones produce large JPEG images (often between 15MB and 20MB). Uploading these raw files over unstable 3G and 4G networks in rural districts like Senapati and Tamenglong frequently resulted in HTTP request timeouts. 

To make the app usable in areas with weak cellular signals, we integrated client-side compression in the browser using the \texttt{browser-image-compression} library. The script downscales images to a maximum width of 1024 pixels before upload, reducing the file size to under 300KB while preserving enough visual detail for the classification model. This client-side optimization resolved both client upload timeouts and the payload size limit in our Next.js proxy route.

The lifecycle of a verification request spans five processing stages, illustrated in Figure~\ref{fig:unified_integration_flow}.

\begin{figure}[htbp]
\centering
\resizebox{\linewidth}{!}{%
\begin{tikzpicture}[
    node distance=1.8cm and 3.5cm,
    box/.style={rectangle, draw=black!80, rounded corners, minimum width=3.8cm, minimum height=1.1cm, align=center, font=\small, thick},
    client/.style={box, fill=blue!10, draw=blue!80},
    proxy/.style={box, fill=yellow!15, draw=yellow!80!black},
    backend/.style={box, fill=orange!15, draw=orange!80!black},
    ai/.style={box, fill=red!15, draw=red!80!black},
    storage/.style={box, fill=purple!15, draw=purple!80!black},
    deploy/.style={box, fill=green!15, draw=green!80!black},
    arrow/.style={->, thick, >=stealth}
]

% Request Flow Nodes (Top Row)
\node[client] (browser) {\textbf{Browser UI}\\(Image Upload)};
\node[proxy, right=of browser] (proxy) {\textbf{Next.js Proxy}\\(Axios / JWT Auth)};
\node[backend, right=of proxy] (backend) {\textbf{ASP.NET Core API}\\(Validation \& Logic)};

% AI & Storage Nodes (Middle Row)
\node[ai, below=of backend] (ai) {\textbf{AI \& ML Services}\\(OCR, ResNet, OpenAI)};
\node[storage, below=of proxy] (db) {\textbf{Cloud Storage}\\(Azure Blob \& SQL DB)};

% Deployment Nodes (Bottom Row)
\node[deploy, below=4.5cm of browser] (docker) {\textbf{Docker Build}\\(Multi-stage CI/CD)};
\node[deploy, right=of docker] (acr) {\textbf{Azure Registry}\\(Container Storage)};
\node[deploy, right=of acr] (aca) {\textbf{Container Apps}\\(Serverless Scaling)};

% Request Flow Connections
\draw[arrow] (browser) -- node[above, font=\footnotesize] {HTTPS / Multipart} (proxy);
\draw[arrow] (proxy) -- node[above, font=\footnotesize] {REST API} (backend);
\draw[arrow] (backend) -- node[right, font=\footnotesize] {Async Inference} (ai);
\draw[arrow] (ai) -- node[above, font=\footnotesize] {Persistence} (db);
\draw[arrow] (db) -- node[left, font=\footnotesize] {Status Response} (proxy);

% Deployment Connections
\draw[arrow, dashed, draw=green!60!black] (docker) -- (acr);
\draw[arrow, dashed, draw=green!60!black] (acr) -- (aca);

% Background grouping boxes
\begin{scope}[on background layer]
    \node[fill=gray!10, rounded corners, fit=(browser) (proxy) (backend) (ai) (db), inner sep=0.4cm] (bg1) {};
    \node[above=0.1cm of bg1.north, font=\small\bfseries, color=gray!70!black] {Operational Verification Pipeline};
    
    \node[fill=green!5, rounded corners, fit=(docker) (acr) (aca), inner sep=0.4cm] (bg2) {};
    \node[above=0.1cm of bg2.north, font=\small\bfseries, color=green!60!black] {Cloud Deployment Pipeline};
\end{scope}

\end{tikzpicture}%
}
\caption{Architectural Workflow and Cloud Deployment Pipeline}
\label{fig:unified_integration_flow}
\end{figure}

The steps in our operational pipeline are structured as follows:

\begin{enumerate}
    \item \textbf{Image Compression and Proxy Routing:} When a farmer uploads an image, the client-side JavaScript compresses the file in the browser. The browser sends the compressed image to our Next.js API route. The proxy server attaches the active JWT authorization token, refreshing it automatically if it has expired \cite{axios2024, jones2015jwt}.
    
    \item \textbf{Backend Ingestion and Initial Validation:} The Next.js proxy forwards the payload to the ASP.NET Core API. The backend processes the request headers and performs validation checks on the image properties before starting the AI analysis \cite{freeman2020pro}.
    
    \item \textbf{Distributed Model Inference:} The backend distributes the image processing load to dedicated microservices. A custom FastAPI service serves our ResNet50 classification model, while Azure Computer Vision handles OCR text extraction. When we first deployed the FastAPI service to serverless container instances on GCP Cloud Run, loading the PyTorch framework and model weights from cold starts caused an unacceptable 15-second latency. To resolve this, we configured our container settings to keep at least one instance active during farming seasons, bringing standard response times down to under two seconds.
    
    \item \textbf{Cloud Storage and Database Persistence:} Once the services complete their runs, the backend stores the results. To prevent database bloat, we avoid storing raw image files as binary data in our SQL database. Instead, images are saved directly to Azure Blob Storage, and the SQL database only stores the resulting storage URLs alongside the model scores and extracted batch data \cite{copard2020azure}.
    
    \item \textbf{Report Synthesis and UI Rendering:} The backend feeds the combined OCR text and seed classification results to Azure OpenAI. The LLM generates a clear, plain-language advisory report. This text is passed back through our API proxy and displayed to the farmer on the web interface.
\end{enumerate}

Alongside this verification pipeline, we run an automated build and release workflow. When code changes are pushed to GitHub, a CI/CD runner compiles our services using multi-stage Dockerfiles. The resulting images are stored in the Azure Container Registry (ACR) and automatically deployed to Azure Container Apps (ACA), which handles horizontal scaling based on request volume \cite{merkel2014docker}.

\subsection{Unexpected Integration Challenges and Resolutions}

During system integration, we encountered several unexpected challenges related to cross-domain security and downstream service latency that required architectural modifications. The first major challenge arose when linking our Next.js frontend, hosted on Vercel, with the ASP.NET Core backend API, hosted on Render. Because these services reside on different domains, modern web browsers blocked authorization cookies due to strict SameSite cookie policies. To resolve this without compromising session security, we routed all client requests through a Next.js server-side API proxy. This server-side route acts as an intermediary, injecting JWT tokens and forwarding request payloads to the backend Render server, successfully bypassing CORS restrictions and cross-domain cookie blocks.

A second integration challenge concerned meeting our non-functional requirement of keeping the end-to-end inference latency under 3.0 seconds. Initially, the backend orchestrated calls to Azure Blob Storage, the ResNet50 FastAPI service, Azure OCR, and Azure OpenAI sequentially. This sequential execution caused cumulative latency to exceed 7.0 seconds per request, resulting in browser timeout errors. To optimize this pipeline, we refactored the ASP.NET Core orchestrator to process independent operations in parallel. By executing the CPU-bound ResNet50 classification and Azure OCR text extraction concurrently using C\# asynchronous tasks (\texttt{Task.WhenAll}), we reduced the overall request-response window to under 2.5 seconds, ensuring a highly responsive user experience even on slower networks.

\subsection{Live Prototype Deployment}
For our live prototype demonstration and evaluation, the Next.js frontend is deployed on Vercel, and the ASP.NET Core backend is hosted on Render. These environments allow us to run real-world tests with active users. The live production endpoints can be accessed at the following URLs:
\begin{itemize}[nosep, topsep=2pt]
    \item Deployed Frontend Client: \href{https://agriverify02.vercel.app/}{\nolinkurl{https://agriverify02.vercel.app/}}
    \item Deployed Backend API Gateway: \href{https://fakeseeddetection.onrender.com/}{\nolinkurl{https://fakeseeddetection.onrender.com/}}
\end{itemize}
The Next.js API proxy routes all client-side queries directly to the Render backend endpoint, which processes the request payloads and coordinates downstream AI inference tasks.